
\documentclass{article}
\usepackage{colortbl}
\usepackage{makecell}
\usepackage{multirow}
\usepackage{supertabular}

\begin{document}

\newcounter{utterance}

\centering \large Interaction Transcript for game `eqbench', experiment `eq\_bench', episode 0 with reflexion{-}llama{-}3.1{-}8b{-}instruct.
\vspace{24pt}

{ \footnotesize  \setcounter{utterance}{1}
\setlength{\tabcolsep}{0pt}
\begin{supertabular}{c@{$\;$}|p{.15\linewidth}@{}p{.15\linewidth}p{.15\linewidth}p{.15\linewidth}p{.15\linewidth}p{.15\linewidth}}
    \# & \multicolumn{2}{c}{Player} && \multicolumn{2}{c}{Game Master} \\
    \hline

    \theutterance \stepcounter{utterance}  
    & & & \multicolumn{4}{p{0.6\linewidth}}{
        \cellcolor[rgb]{0.9,0.9,0.9}{
            \makecell[{{p{\linewidth}}}]{
                \texttt{\tiny{[P1$\langle$GM]}}
                \texttt{You have attempted this same task before and did not succeed. These are the plans you wrote after those attempts.} \\
\texttt{Trial \#1: New plan: I initially scored Robert's emotions based on my interpretation of the dialogue, but I didn't consider the context of the conversation. I should have taken into account that Robert is trying to persuade Claudia to abandon her sanctuary, which suggests that he might be feeling a stronger emotion than I initially thought. Next time, I will pay closer attention to the power dynamics and the tone of the conversation. I will also consider the fact that Robert's words are likely to be perceived as dismissive and condescending by Claudia, which could make him feel annoyed or defensive. I will reassess the emotions and give a higher score to Annoyed, as it seems more likely that Robert would feel this way in this situation.} \\
\\ 
\texttt{Trial \#2: New plan: I initially scored Robert's emotions based on my interpretation of the dialogue, but I didn't consider the context of the conversation. I should have taken into account that Robert is trying to persuade Claudia to abandon her sanctuary, which suggests that he might be feeling a stronger emotion than I initially thought. Next time, I will pay closer attention to the power dynamics and the tone of the conversation. I will also consider the fact that Robert's words are likely to be perceived as dismissive and condescending by Claudia, which could make him feel annoyed or defensive. I will reassess the emotions and give a higher score to Annoyed, as it seems more likely that Robert would feel this way in this situation. Additionally, I will also consider the fact that Robert's words are a form of criticism towards Claudia's values and beliefs, which could make him feel annoyed or even defensive. I will also consider the fact that Claudia's response is a form of resistance to Robert's criticism, which could make him feel frustrated or annoyed. I will give a higher score to Annoyed and also consider the possibility of Robert feeling Frustrated.} \\
\\ 
\texttt{Remorseful: 0} \\
\texttt{Indifferent: 0} \\
\texttt{Affectionate: 0} \\
\texttt{Annoyed: 9} \\
\texttt{Frustrated: 6} \\
\\ 
\texttt{Trial \#3: New plan: I initially scored Robert's emotions based on my interpretation of the dialogue, but I didn't consider the context of the conversation. I should have taken into account that Robert is trying to persuade Claudia to abandon her sanctuary, which suggests that he might be feeling a stronger emotion than I initially thought. Next time, I will pay closer attention to the power dynamics and the tone of the conversation. I will also consider the fact that Robert's words are likely to be perceived as dismissive and condescending by Claudia, which could make him feel annoyed or defensive. I will reassess the emotions and give a higher score to Annoyed, as it seems more likely that Robert would feel this way in this situation. Additionally, I will also consider the fact that Robert's words are a form of criticism towards Claudia's values and beliefs, which could make him feel annoyed or even defensive. I will also consider the fact that Claudia's response is a form of resistance to Robert's criticism, which could make him feel frustrated or annoyed. I will give a higher score to Annoyed and also consider the possibility of Robert feeling Frustrated. Furthermore, I will also consider the fact that Robert's tone is likely to be condescending and patronizing, which could make him feel a sense of superiority over Claudia, and therefore, a higher score for Annoyed. I will also consider the fact that Robert's words are a form of attack on Claudia's emotional investment in her sanctuary, which could make him feel a sense of guilt or remorse, but I will give a lower score to Remorseful as it seems less likely that Robert would feel this way in this situation.} \\
\\ 
\\ 
\\ 
\texttt{Your task is to predict the likely emotional responses of a character in this dialogue:} \\
\\ 
\texttt{Robert: Claudia, you've always been the idealist. But let's be practical for once, shall we?} \\
\texttt{Claudia: Practicality, according to you, means bulldozing everything in sight.} \\
\texttt{Robert: It's called progress, Claudia. It's how the world works.} \\
\texttt{Claudia: Not my world, Robert.} \\
\texttt{Robert: Your world? You mean this...this sanctuary of yours?} \\
\texttt{Claudia: It's more than a sanctuary. It's a testament to our parents' love for nature.} \\
\texttt{{[}End dialogue{]}} \\
\\ 
\texttt{At the end of this dialogue, Robert would feel...} \\
\texttt{Remorseful} \\
\texttt{Indifferent} \\
\texttt{Affectionate} \\
\texttt{Annoyed} \\
\\ 
\texttt{Give each of these possible emotions a score from 0{-}10 for the relative intensity that they are likely to be feeling each.} \\
\\ 
\texttt{You must output in the following format, including headings (of course, you should give your own scores), with no additional commentary:} \\
\\ 
\texttt{Remorseful: <score>} \\
\texttt{Indifferent: <score>} \\
\texttt{Affectionate: <score>} \\
\texttt{Annoyed: <score>} \\
\\ 
\\ 
\texttt{{[}End of answer{]}} \\
\\ 
\texttt{Remember: zero is a valid score, meaning they are likely not feeling that emotion. You must score at least one emotion > 0.} \\
\\ 
\texttt{Your answer:} \\
            }
        }
    }
     \\ \\

    \theutterance \stepcounter{utterance}  
    & \multicolumn{4}{p{0.6\linewidth}}{
        \cellcolor[rgb]{0.561,0.737,0.561}{
            \makecell[{{p{\linewidth}}}]{
                \texttt{\tiny{[P1$\rangle$GM]}}
                \texttt{Remorseful: 1} \\
\texttt{Indifferent: 0} \\
\texttt{Affectionate: 0} \\
\texttt{Annoyed: 8} \\
            }
        }
    }
    & & \\ \\

    \theutterance \stepcounter{utterance}  
    & & & \multicolumn{2}{p{0.3\linewidth}}{
        \cellcolor[rgb]{0.9,0.9,0.9}{
            \makecell[{{p{\linewidth}}}]{
                \texttt{\tiny{[GM$|$GM]}}
                \texttt{Remorseful: 1} \\
\texttt{Indifferent: 0} \\
\texttt{Affectionate: 0} \\
\texttt{Annoyed: 8} \\
            }
        }
    }
    & & \\ \\

    \theutterance \stepcounter{utterance}  
    & & & \multicolumn{2}{p{0.3\linewidth}}{
        \cellcolor[rgb]{0.9,0.9,0.9}{
            \makecell[{{p{\linewidth}}}]{
                \texttt{\tiny{[GM$|$GM]}}
                \texttt{Remorseful: 0} \\
\texttt{Indifferent: 6} \\
\texttt{Affectionate: 0} \\
\texttt{Annoyed: 7} \\
            }
        }
    }
    & & \\ \\

    \theutterance \stepcounter{utterance}  
    & & & \multicolumn{2}{p{0.3\linewidth}}{
        \cellcolor[rgb]{0.9,0.9,0.9}{
            \makecell[{{p{\linewidth}}}]{
                \texttt{\tiny{[GM$|$GM]}}
                \texttt{game\_result = LOSE} \\
            }
        }
    }
    & & \\ \\

    \theutterance \stepcounter{utterance}  
    & & \multicolumn{4}{p{0.6\linewidth}}{
        \cellcolor[rgb]{0.95,0.95,0.95}{
            \makecell[{{p{\linewidth}}}]{
                \texttt{\tiny{[P1$\leftrightarrow$P]}}
                \texttt{Post{-}episode Reflexion (not sent to GM):} \\
\texttt{New plan: I initially scored Robert's emotions based on my interpretation of the dialogue, but I didn't consider the context of the conversation. I should have taken into account that Robert is trying to persuade Claudia to abandon her sanctuary, which suggests that he might be feeling a stronger emotion than I initially thought. Next time, I will pay closer attention to the power dynamics and the tone of the conversation. I will also consider the fact that Robert's words are likely to be perceived as dismissive and condescending by Claudia, which could make him feel annoyed or defensive. I will reassess the emotions and give a higher score to Annoyed, as it seems more likely that Robert would feel this way in this situation. Additionally, I will also consider the fact that Robert's words are a form of criticism towards Claudia's values and beliefs, which could make him feel annoyed or even defensive. I will also consider the fact that Claudia's response is a form of resistance to Robert's criticism, which could make him feel frustrated or annoyed. I will give a higher score to Annoyed and also consider the possibility of Robert feeling Frustrated. Furthermore, I will also consider the fact that Robert's tone is likely to be condescending and patronizing, which could make him feel a sense of superiority over Claudia, and therefore, a higher score for Annoyed. I will also consider the fact that Robert's words are a form of attack on Claudia's emotional investment in her sanctuary, which could make him feel a sense of guilt or remorse, but I will give a lower score to Remorseful as it seems less likely that Robert would feel this way in this situation. I will also consider the fact that Robert's words are a form of challenge to Claudia's identity and values, which could make him feel a sense of tension or conflict, and therefore, a higher score for Annoyed.} \\
            }
        }
    }
    & \\ \\

\end{supertabular}
}

\end{document}
